\section{Growing-dimensional extension and a metric-free inversion baseline}
\label{sec:scalable}
The three-level mechanism is not a fixed-dimensional accident.  We now embed the same
eigenvector-coalescence geometry into a genuinely coupled family of increasing dimension.
The purpose is mathematical: to show that the separation between metric conditioning and
Euclidean inversion stability persists uniformly with dimension.  A standard Hermitian
dilation is then recorded only as a computational implication: it uses singular values rather
than a pseudo-Hermitian metric and therefore makes explicit why metric ill-conditioning does
not, by itself, imply hard spectral inversion.

\subsection{A coupled bidiagonal family with bounded inverse}
\begin{theorem}[Nonzero spectral coalescence in every dimension]
\label{thm:growingchain}
For an integer $N\geq3$ and $0<\varepsilon\leq1/2$, let
\begin{equation}\label{eq:growingchain}
 \begin{aligned}
 C_{N,\varepsilon}&=\operatorname{diag}(d_1,\ldots,d_N)
  +\sum_{j=1}^{N-1}e_je_{j+1}^{\mathsf T},\\
 d_1&=1,\quad d_2=1+\varepsilon,\quad
 d_j=2+\frac{j-3}{N-2}\quad (3\leq j\leq N).
 \end{aligned}
\end{equation}
The matrix has $N$ distinct, positive eigenvalues, is diagonalizable and
admits positive compatible metrics. Nevertheless,
\begin{equation}\label{eq:growing_bounds}
\kappa_*(C_{N,\varepsilon})\geq\frac{8}{\varepsilon^2},
\qquad \|C_{N,\varepsilon}\|_2\leq4,
\qquad \|C_{N,\varepsilon}^{-1}\|_2\leq\sqrt6,
\qquad r_{\mathrm{sing}}(C_{N,\varepsilon})\geq\frac1{\sqrt6}.
\end{equation}
Thus the best metric-based Euclidean invertibility certificate is at most
$\varepsilon/(2\sqrt2)$ while the exact distance to singularity remains uniformly
positive in \emph{both} $N$ and $\varepsilon$. Moreover, if a scalar polynomial
$q$ satisfies $q(C_{N,\varepsilon})=C_{N,\varepsilon}^{-1}$,
then $\deg q\geq N-1$.
\end{theorem}
\begin{proof}
The diagonal entries are distinct and lie in $[1,3)$, so the upper triangular
matrix is diagonalizable with positive spectrum. The first three right eigenvectors are supported on
$S=\operatorname{span}\{e_1,e_2,e_3\}$ and coincide with those of
$A_\varepsilon$ in~\eqref{eq:threelevel}.  Order an eigenvector matrix so that
\begin{equation}
 V=\begin{pmatrix}V_\varepsilon&R\\0&V_{\rm tail}\end{pmatrix},
\end{equation}
where $V_\varepsilon$ is the three-dimensional matrix in
\eqref{eq:chain_V}.  For any admissible metric,
$\eta=V^{-\dagger}WV^{-1}$ with a positive diagonal $W$ (the eigenvalues
of $C_{N,\varepsilon}$ are simple).  For the vector
$x=2^{-1/2}(0,1,-1,0,\ldots,0)^{\mathsf T}$, block triangularity gives
$V^{-1}x=(V_\varepsilon^{-1}x_{1:3},0)^{\mathsf T}$, so
$x^\dagger\eta x$ is exactly the same three-weight expression used in
\eqref{eq:xetax}.  For $y=e_1$, the identity
$\eta^{-1}=VW^{-1}V^\dagger$ shows that $y^\dagger\eta^{-1}y$ contains the
three positive contributions appearing in~\eqref{eq:yetainvy}, plus
nonnegative contributions from the remaining eigendirections.  Hence the
Rayleigh-product argument of Theorem~\ref{thm:metric_sandwich} applies
without replacing the inverse of a compression by a compressed inverse and yields
$\kappa_2(\eta)\geq8/\varepsilon^2$.  Taking the infimum gives
$\kappa_*(C_{N,\varepsilon})\geq8/\varepsilon^2$.
The strictly upper shift has spectral norm one and $\max_jd_j<3$,
whence $\|C_{N,\varepsilon}\|_2\leq4$.
The inverse is upper triangular with
\begin{equation}\label{eq:growinginverse}
 (C_{N,\varepsilon}^{-1})_{ij}=
 \begin{cases}(-1)^{j-i}\displaystyle\prod_{k=i}^{j}d_k^{-1},&i\leq j,\\
 0,&i>j.\end{cases}
\end{equation}
For rows $i\geq3$, the sum of absolute entries is at most
$\sum_{r\geq1}2^{-r}=1$; the corresponding sums for rows two and one
are at most two and three. Hence $\|C^{-1}\|_\infty\leq3$.
For columns $j\geq3$, the entries in rows $i\geq3$ sum to at most one,
while those in rows one and two sum to at most one together.
Columns one and two have absolute sums at most one and two, respectively.
Thus $\|C^{-1}\|_1\leq2$, and
$\|C^{-1}\|_2\leq\sqrt{\|C^{-1}\|_1\|C^{-1}\|_\infty}\leq\sqrt6$.
Equation~\eqref{eq:radius} gives the singularity-radius bound.
Finally, if $\deg q\leq N-2$, then $xq(x)-1$ has degree at most
$N-1$ but vanishes at the $N$ distinct eigenvalues of $C$.
It would therefore vanish identically, contradicting its constant
term $-1$. Hence $\deg q\geq N-1$.
\end{proof}
\begin{remark}[What the polynomial degree does and does not count]
The degree lower bound is \emph{only} for exact scalar polynomials in
$C_{N,\varepsilon}$. A monomial-product and linear-combination implementation
of such an exact polynomial has a branch using at least $N-1$ input-matrix
blocks. The statement is \emph{not} a query lower bound for an arbitrary
quantum linear-system algorithm, which may use $C^\dagger$, a Hermitian
dilation, rational approximations, or singular-value transformations.
In particular, $N=2^q$ would turn this particular exact polynomial
implementation into one with exponential degree in the data-qubit count $q$.
No gate count for an oracle encoding this coupled chain is assumed here.
\end{remark}

\subsection{A standard Hermitian-dilation baseline}
\begin{proposition}[A metric-free Hermitian reduction]
\label{prop:hermitian_dilation}
Let $A$ be invertible with a unitary block encoding
$\langle0^a|U_A|0^a\rangle=A/\alpha$.
On an additional flag qubit define
\begin{equation}\label{eq:herm_dilation}
 \mathcal H_A=\begin{pmatrix}0&A\\A^\dagger&0\end{pmatrix},\qquad
 \mathcal U_A=|0\rangle\langle1|\otimes U_A
 +|1\rangle\langle0|\otimes U_A^\dagger.
\end{equation}
Then $\mathcal U_A$ is unitary and, on the all-zero original
block-encoding ancilla, encodes $\mathcal H_A/\alpha$.
The nonzero spectral magnitudes of $\mathcal H_A$ are exactly the singular
values of $A$, and
\begin{equation}\label{eq:dilationbaseline}
 \kappa_2(\mathcal H_A)=\kappa_2(A),\qquad
 \mathcal H_A^{-1}(|0\rangle\otimes b)
 =|1\rangle\otimes A^{-1}b.
\end{equation}
For the growing coupled chain, $\kappa_2(\mathcal H_C)\leq4\sqrt6$,
independently of $N$ and $\varepsilon$.
\end{proposition}
\begin{proof}
The two off-diagonal blocks of $\mathcal U_A$ are unitaries which
exchange orthogonal flag sectors, making $\mathcal U_A$ unitary.
Projecting the old ancilla gives the advertised block.
Squaring $\mathcal H_A$ gives
$\operatorname{diag}(AA^\dagger,A^\dagger A)$;
its eigenvalue magnitudes are the singular values of $A$.
Direct block inversion gives~\eqref{eq:dilationbaseline}.
Apply Theorem~\ref{thm:growingchain} to obtain the last bound.
\end{proof}
\begin{remark}
Hermitian dilation and singular-value-based linear solvers are established
techniques~\cite{Gilyen2019,Berry2026}.  Proposition~\ref{prop:hermitian_dilation}
is included as a diagnostic consequence rather than an algorithmic novelty claim.  It shows
that the diverging optimal metric condition number is not inherited by the singular-value
conditioning of the dilation.  Any end-to-end quantum cost still depends on the access model,
state preparation and readout.
\end{remark}
